\SourceLeaf{471}
\par\Needspace{12\baselineskip}\bigskip
\pdfbookmark[2]{Exposition, Fourth section, First draft}{pallas-II-section4-first}
\begin{center}
{\large\bfseries Fourth section.\quad [First draft.]\par}
\medskip
\textit{Principles for determining the motion for an unlimited time.}
\end{center}
\pdfbookmark[3]{Exposition, [21]}{pallas-II-article21}
\begin{center}[21.]\end{center}
\noindent\hspace*{\parindent}Integrating the instantaneous variations of the elements from their numerical values gives the elements for only a limited time. To carry out the indefinite integration, one must first express these variations as functions of time that can be integrated. This is precisely where the difficulty of the problem lies. Given the smallness of the perturbing masses compared with the mass of the Sun, we regard the variations of the elements, or more generally the quantities containing the factor \(\mu\), as very small quantities of the \textit{first order}. Those containing its square or the product of \(\mu\) and another perturbing mass will be of the second order, and so on. From now on we shall denote the true longitude of the node (reckoned on the fixed ecliptic) at an undetermined instant by \(\AscendingNode+\delta\AscendingNode\), where \(\AscendingNode\) is a constant quantity and \(\delta\AscendingNode\) is the variable part of the first order. We shall proceed in the same way with the other elements. It is evident from this that if we determine the instantaneous variation of an element by calculating all the quantities in its expression from the constant parts of the elements, that variation will be accurate up to quantities of the second order (excluded). In general, the error in the variations of the elements will be one order higher than the order of the values of the elements used to calculate them. This gives rise to the method of successive approximations. We must admit that, however simple this method is in principle, carrying it out could become almost impracticable if even the second approximation had to be calculated \textit{in full}. Fortunately this is quite unnecessary for the needs of astronomy. We shall confine ourselves to the first approximation and add a few modifications derived from a reasoned consideration of the second-order terms.
